\section{Study Area, Data, and Forecast Protocol}
\label{sec:data}
% Numbers in this section come from reports/section3_facts.md (generated by scripts/section3/facts.py); section numbers of that file are given in the comments.

\subsection{Study Area and River Network}
\label{sec:study}
We study the Tennessee River Basin (hydrologic region 06), which drains about 106,000~km$^2$ of the southeastern United States (Figure~\ref{fig:area}a). Mean catchment elevation rises from about 190~m in the western lowlands (west of 87$^\circ$W) to about 700~m in the east (east of 84$^\circ$W), and reaches more than 1,500~m in the highest catchments of the Appalachian highlands. Mean annual precipitation is about 1,420~mm (AORC, 2000--2019). It differs little between the western and eastern parts of the basin (about 1,440 and 1,360~mm), and only the highest catchments are clearly wetter (about 1,590~mm).
% facts sheet section 15 (precipitation by region and elevation class); section 14 (AORC units assumed as mm per six-hour step).

We represent the river network by the flowlines of NHDPlus Version 2.1 \todocite{NHDPlus}. The network has 58,834 reaches connected by 60,267 directed links, with Strahler stream orders from 1 to 8. Most reaches drain small areas: 45,703 reaches (77.7\%) drain less than 50~km$^2$, 7,868 drain 50--500~km$^2$, 3,694 drain 500--5,000~km$^2$, and 1,569 drain more than 5,000~km$^2$. Every reach is a forecast location.

\begin{figure}
\figfile{figures/fig1_study_area.pdf}{0.85\textwidth}
\caption{Study area. (a) River network of the Tennessee River Basin (NHDPlus Version 2.1; 58,834 reaches, major rivers drawn) and the 101 evaluation gauges. (b) The three basins withheld in the blocked protocol, B1--B3, with the gauges evaluated in each; the other evaluation gauges are shown in gray.}
\label{fig:area}
\end{figure}

\subsection{Streamflow Observations}
\label{sec:gauges}
Observed discharge comes from two public sources. Until 2019-10-31 we use the hourly streamflow of the CAMELS-H data set \todocite{CAMELSH}, and from 2019-11-01 we use instantaneous discharge from the USGS National Water Information System (NWIS), averaged to hourly values. Hourly values are averaged to six-hour means over $[t, t+6\,\mathrm{h})$ in UTC, and a six-hour mean is formed only when all six hours are present. Negative values are treated as missing. No other quality control is applied and gaps are not filled.

We matched 128 USGS gauges to reaches of the network: each gauge was assigned the nearest reach within 1,500~m whose drainage area is between 0.8 and 1.25 times the gauge's reported area. Of these gauges, 125 are CAMELS-H sites and 3 have NWIS records only, beginning in 2019. We evaluate forecasts at the CAMELS-H gauges that have at least 100 forecast--observation pairs at 24~h lead in the test period and a positive mean flow with non-negligible variability (standard deviation of at least $10^{-3}$~m$^3$\,s$^{-1}$). This gives 101 evaluation gauges. Of the other 27 gauges, 24 have records that end before the test period, and 3 have no CAMELS-H record; gauges without a CAMELS-H record are not scored.

The evaluation gauges are not a representative sample of the network (Figure~\ref{fig:gauges}a). Their median drainage area is 382~km$^2$, against 5.8~km$^2$ for all reaches, and only 4.0\% of them drain less than 50~km$^2$, against 77.7\% of the reaches. Their drainage areas range from 14 to 5,306~km$^2$. Skill at gauges therefore describes medium and large rivers. We return to the small streams with the reach-level evaluation in Section~\ref{sec:holdout}.

Because RQ3 concerns observations made upstream of a forecast location, we distinguish two kinds of evaluation gauge by their position in the network (Figure~\ref{fig:gauges}b). A gauge is \emph{nested} if at least one of the 128 matched gauges lies in its upstream basin, and a \emph{headwater gauge} otherwise. Of the 101 evaluation gauges, 44 are nested and 57 are headwater gauges, with median drainage areas of 1,238 and 275~km$^2$. Nested gauges therefore drain larger basins, and we compare results by position only at similar drainage areas.
% facts sheet section 15. If only the 101 evaluation gauges count as upstream gauges, the split is 39 nested (median 1,659 km^2) and 62 (median 276 km^2).

Record length varies among the evaluation gauges (Figure~\ref{fig:gauges}c). Over the training years (2000--2018), 19 gauges are observed less than half of the time and 3 gauges (USGS 03566420, 03572110, and 03591800) not at all, because their records begin in 2019 or later. During the test period, most gauges are observed almost continuously: 89 of the 101 are observed at least 95\% of the time.

\begin{figure}
\figfile{figures/fig2_gauges.pdf}{\textwidth}
\caption{Evaluation gauges. (a) Cumulative distribution of drainage area for all reaches ($n = 58{,}834$), all matched gauges ($n = 128$, dashed) and evaluation gauges ($n = 101$); dots mark the medians. (b) Drainage area of headwater ($n = 57$) and nested ($n = 44$) evaluation gauges; bars mark the medians. (c) Fraction of six-hour values observed at each evaluation gauge in the training and test periods, each sorted separately.}
\label{fig:gauges}
\end{figure}

\subsection{Meteorological Forcing and Catchment Attributes}
\label{sec:inputs}
Weather before the forecast issue time comes from the Analysis of Record for Calibration (AORC), version 1.1 \todocite{AORC}. We use precipitation, 2~m air temperature and specific humidity, surface pressure, downward shortwave and longwave radiation, and the two components of the 10~m wind. The hourly fields on the 1/120$^\circ$ grid are aggregated to six-hour values and averaged over the local catchment of each reach, weighted by area. The 1,192 reaches without a catchment polygon in NHDPlus receive no local forcing.

Precipitation after the issue time comes from the control member of the Global Ensemble Forecast System, version 12 \cite{Zhou2022,Guan2022}. We use six-hour totals from the 00Z cycle at leads to 240~h on the native 0.25$^\circ$ grid. Each AORC cell takes the value of the nearest forecast grid point, and the values are averaged over catchments with the same weights as the observed forcing. Training and validation use the GEFSv12 reforecast \cite{Guan2022}, and testing uses the archived operational forecasts. No other variable is forecast, so future temperature, humidity, radiation, and wind are unknown to all models.

Each reach is described by 24 static attributes: reach length and slope, upstream and downstream elevation, stream order, local and upstream drainage area, upstream channel length from NHDPlus, and catchment and watershed means of mean elevation and of seven soil and subsurface properties (clay, sand, and organic matter content, permeability, depth to bedrock, depth to the water table, and hydraulic conductivity) from EPA StreamCat \todocite{StreamCat}. From the AORC record we also derive seven antecedent-wetness indices for each reach: exponentially weighted precipitation with time scales from 7 to 365 days and air temperature with time scales of 30 and 90 days.
% NOTE: list the seven indices and their time scales in Table S1 (processed/hydro/antecedent.json).

\subsection{National Water Model}
\label{sec:nwm}
The National Water Model (NWM) is used in two roles. Its version 3.0 retrospective simulation \cite{Cosgrove2024} gives continuous discharge at every reach, which we use as a pretraining target (Section~4) and as the reference for the reach-level evaluation. The archived operational medium-range forecast (deterministic member, 00Z cycle, flow every 3~h to 240~h) is the reference forecast at gauges. The operational model changed during the test period, from version 2.0 to 2.1 on 2021-04-20 and to 3.0 on 2023-09-20. The operational NWM nudges simulated flow toward observations at gauges before each forecast \cite{Cosgrove2024}, including gauges that our models treat as withheld, so it serves as a reference and not as a competitor under the same information.
% NOTE for Section 4: the LSTM baseline is trained without NWM pretraining and without observations of other gauges; state this where RQ1 is set up.

Table~\ref{tab:data} summarizes the data products and their periods.

\begin{table}
\caption{Data Products}
\label{tab:data}
\footnotesize
\setlength{\parindent}{0pt}
\renewcommand{\arraystretch}{1.3}
\setlength{\tabcolsep}{3pt}
\begin{tabular}{@{}>{\raggedright\arraybackslash}p{0.23\textwidth} >{\raggedright\arraybackslash}p{0.22\textwidth} >{\raggedright\arraybackslash}p{0.22\textwidth} >{\raggedright\arraybackslash}p{0.23\textwidth}@{}}
\toprule
\textbf{Product} & \textbf{Use} & \textbf{Period} & \textbf{Resolution} \\
\midrule
AORC v1.1 & Observed weather & 1979-02-01 to 2025-12-31 & Hourly, 1/120$^\circ$ \\
\addlinespace
GEFSv12 reforecast, control member & Forecast precipitation, training and validation & Issued 2000-01-01 to 2019-12-31 & Six-hour totals to 240~h, 0.25$^\circ$ \\
\addlinespace
GEFSv12 operational archive, control member & Forecast precipitation, testing & Issued 2020-09-23 to 2025-12-31 & Six-hour totals to 240~h, 0.25$^\circ$ \\
\addlinespace
CAMELS-H & Observed discharge & 1984-10-12 to 2019-10-31 & Hourly, 125 gauges \\
\addlinespace
USGS NWIS & Observed discharge & 2019-11-01 to 2025-12-31 & Instantaneous, 103 gauges with data \\
\addlinespace
NHDPlus V2.1, EPA StreamCat & River network, static attributes & --- & 58,834 reaches \\
\addlinespace
NWM v3.0 retrospective & Pretraining target, reach-level reference & 1979-02-01 to 2023-01-31 & Hourly, 58,834 reaches \\
\addlinespace
NWM operational medium-range forecast & Reference forecast at gauges & Issued 2020-09-23 to 2025-12-22 & Every 3~h to 240~h, deterministic member \\
\bottomrule
\end{tabular}
\begin{flushleft}
\footnotesize \emph{Note.} All products are aggregated to six-hour values. Dates give the extent of each record over the gauges and reaches used; the NWM retrospective is given as published, and the reach-level evaluation uses its targets up to 2023-01-30.
\end{flushleft}
\end{table}

\subsection{Forecast Setup and Temporal Partitions}
\label{sec:protocol}
A forecast is issued every day at 00Z and gives six-hour mean discharge at every reach for 40 lead times, from 6 to 240~h. At the issue time, a model has the observed weather of the previous 36 days at every reach, forecast precipitation for the next 10 days, and the discharge observed up to the issue time at gauges that are not withheld (Figure~\ref{fig:protocol}). A gauge contributes observations only if it reported at least 8 of the last 32 six-hour values. We assume that observations are available without delay; in real time, NWIS data arrive with a lag, so this assumption favors the use of upstream observations.

The data are split in time. Training uses the 6,931 reforecast issue days from 2000-01-01 to 2018-12-22, so that all training targets precede 2019. Validation uses the 355 issue days from 2019-01-01 to 2019-12-21 and serves only for model selection. The test period contains the 1,917 issue days from 2020-09-23, the start of the operational archive, to 2025-12-22 on which both GEFSv12 and NWM forecasts are available. The months between validation and test are not used. Unless stated otherwise, we summarize skill over the leads of 24, 48, 72, 96, and 120~h (days 1--5); the metrics are defined in Section~4.

\begin{figure}
\figfile{figures/fig3_protocol.pdf}{0.85\textwidth}
\caption{Information available at the forecast issue time (00Z). Observed weather at every reach and observed discharge at gauges that are not withheld are available for the preceding 36 days; forecast precipitation is available for the following 240~h. Skill is summarized over leads of 24--120~h (days 1--5).}
\label{fig:protocol}
\end{figure}

\subsection{Spatial Evaluation Protocols}
\label{sec:holdout}
Forecasts at unmonitored reaches cannot be verified directly, so we evaluate them at gauges whose observations are withheld from the model. Three protocols withhold progressively more information, and a fourth evaluation covers every reach.

\textit{Gauged.} Models are trained with all gauges that have training-period records and are evaluated at the 101 evaluation gauges. This conventional setting serves as an upper reference. The three gauges without training-period records are effectively unseen even in this protocol.

\textit{Random hold-out.} The 122 gauges with a CAMELS-H record before 2019 are divided into five folds of 24 or 25 gauges. The gauges are sorted by drainage area, and each consecutive group of five is assigned a random permutation of the fold numbers, so that every fold contains small and large basins. For each fold, a separate model is trained and selected without the observations of that fold's gauges and is evaluated at them. Ninety-eight of the 101 evaluation gauges are evaluated in this protocol; the three gauges without training-period records belong to no fold. A withheld gauge usually has gauged neighbors upstream, downstream, or in adjacent tributaries.

\textit{Blocked basins.} To remove nearby gauged information, we withhold entire upstream basins. Starting with the gauges whose basins contain the most CAMELS-H gauges with training-period records, we select up to three non-overlapping basins that each contain 8 to 30 such gauges and no more than 25\% of the network area. The rule uses only network topology and gauge counts. It selects the basins of the French Broad River near Newport, Tennessee (B1; USGS 03455000), the Holston River near Knoxville, Tennessee (B2; 03495500), and the Hiwassee River at Charleston, Tennessee (B3; 03566000) (Figure~\ref{fig:area}b and Table~\ref{tab:blocked}). They contain 15, 13, and 12 such gauges and cover 4.6, 9.1, and 5.7\% of the network area. For each basin, a model is trained and selected without the observations of all gauges inside it and is evaluated at the 12, 11, and 6 of them that are evaluation gauges, 29 in total. Gauges downstream of a withheld basin remain in training.

The three blocked basins lie in the eastern highlands and differ from the region of the other evaluation gauges (Table~\ref{tab:blocked}). Their gauged catchments are higher (691 to 940~m against 623~m) and more forested (66 to 81\% against 64\%). B1 is also steeper, B2 is drier, and B1 and B3 are wetter than the other gauged catchments. The blocked protocol therefore tests transfer to unseen terrain as well as to unseen basins. With three basins, its results are informative but uncertain.

In all protocols, a withheld gauge contributes no observations at any time, neither as a training target, nor as an input, nor for model selection. We verified this with tests of the data pipeline and of the trained models (Text~S1).

\begin{table}
\caption{Blocked Basins and the Other Evaluation Gauges}
\label{tab:blocked}
\footnotesize
\setlength{\parindent}{0pt}
\renewcommand{\arraystretch}{1.3}
\setlength{\tabcolsep}{4pt}
\begin{tabular}{@{}lrrrrrrr@{}}
\toprule
 & Area & & Evaluated & Elevation & Slope & Precipitation & Forest \\
 & (km$^2$) & Reaches & gauges & (m) & (m m$^{-1}$) & (mm yr$^{-1}$) & (\%) \\
\midrule
B1 & 4,801 & 2,847 & 12 & 940 & 0.070 & 1,596 & 81 \\
B2 & 9,639 & 5,235 & 11 & 774 & 0.039 & 1,247 & 66 \\
B3 & 5,971 & 3,228 & 6 & 691 & 0.048 & 1,621 & 80 \\
Other gauges & --- & --- & 72 & 623 & 0.041 & 1,505 & 64 \\
\bottomrule
\end{tabular}
\begin{flushleft}
\footnotesize \emph{Note.} Area and reaches refer to the whole blocked basin. The other columns are means over the evaluated gauges of values for each gauge's upstream basin: mean elevation (EPA StreamCat), area-weighted reach slope (NHDPlus), mean annual precipitation (AORC, 2000--2019), and forest cover (NLCD 2011).
\end{flushleft}
\end{table}
% facts sheet section 14. Precipitation units (mm per six-hour step) are assumed from the repository's data notes; the per-gauge annual totals (1,160-1,970 mm) are the check. Slope is a reach slope, not a catchment slope.

\textit{Every reach.} Gauges cover few small streams, so we also evaluate forecasts at every reach against the NWM retrospective simulation, for every second test issue up to 2023-01-20 (425 issues). Reaches inside the upstream basin of a randomly withheld gauge (31,684 reaches) are evaluated with the model of that gauge's fold. The remaining reaches lie outside every gauge basin (25,588) or inside the basin of a gauge that belongs to no fold (1,562); they are evaluated with the model trained with all gauges. Because the reference is itself a model, this evaluation measures agreement with the NWM simulation rather than accuracy against observations, and it is used to examine how skill changes with drainage area.
